\section{Frozen inputs, local orders, and authority bounds}
\label{app:local}

\subsection{Exact algebraic input}

The primitive coefficient array of \cref{eq:frozen-cubic} has $20$ entries.
The four Grassmann-chart equations are obtained by the direct substitution
\[
 (u_0,u_1,u_2,u_3)=(s,t,as+ct,bs+dt)
\]
and equating the coefficients of $s^3,s^2t,st^2,t^3$.  A degree-order basis
has Hilbert counts $(1,4,10,12,0)$; lexicographic conversion produces an
eliminant $g(d)$ of degree $27$ and rational back-substitutions for
$a,b,c$.  Substituting those expressions into all four chart equations,
clearing denominators, and reducing modulo $g$ gives four zero remainders.
The complement of the chart is empty by five independent unit-ideal
calculations.

The maximal-order calculation uses a transformed monic degree-$27$
polynomial together with an oriented isomorphism back to the original class
of $d$.  The orientation is part of the input: the original eliminant
generator maps to a rational polynomial in the transformed generator.  The
full numerator vector is too large for print, so the following canonical
guards identify it and the integral basis without abbreviation at the data
layer.

\begin{table}[ht]
\centering
\caption{Compact guards for the exact field presentation.}
\label{tab:field-input-guards}
\begin{tabularx}{\textwidth}{@{}l X@{}}
\toprule
object&SHA-256 or exact property\\
\midrule
oriented original-generator image&
\hexsha{aea67936d6dd88a7ef9ecaabd4740b737e52de9c34958e77bcefc5a51eed0e23}\\
transformed monic polynomial&
\hexsha{1ecdb17baaa586a5361cfdf71ae6af218e05adb0b784aba514410a933fc74c64}\\
canonical $27$-element integral basis&
\hexsha{5102fd87524498ef170b944f174b9c917279f7498f59b372de85b9913a0fbcc8}\\
PARI textual basis carrier&
\hexsha{ec66da41c6ad35e441c9916d3e448206ecc36e857fbab814475420a25c5d4813};
$4{,}549{,}955$ bytes\\
maximality certification&
no unresolved prime returned by \texttt{nfcertify}\\
\bottomrule
\end{tabularx}
\end{table}

The semantics of \texttt{nfbasis}, \texttt{nfinit},
\texttt{idealprimedec}, \texttt{nfcertify}, and \texttt{polsturm} are those
of the PARI/GP user's guide \citep{parigroup2021}.  The software reference
does not certify any displayed output; exact carriers and independent
reconstruction do.

\subsection{Divided discriminant engines}

The Macaulay construction uses a $56\times56$ matrix and a $24\times24$
extraneous block.  Fraction-free determinant evaluation and an independent
exact polynomial engine give the same quotient.  Its decimal-newline guard
is
\[
\texttt{2be931285c05779d59f6d7f9f7006bd87a3c36e72c4d57dab7315413a3672ed9},
\]
and its prime factorization is exactly
\cref{eq:surface-divided-discriminant}.  The field discriminant, computed
from the integral basis rather than this resultant, is positive with $586$
digits and guard
\[
\texttt{7548db5eb3f1c5549d80f6125521e9f3c7f965fb39b7198a8d28f93e8f78d6ca}.
\]
The different origin of these two integers is why the nine-prime surface
envelope and eight-prime field support remain separate throughout the paper.

\subsection{Complete local-row summary}

For readability, repeated rows are compressed only by their multiplicity.
\begin{longtable}{@{}c r r r r r r@{}}
\caption{All maximal-order local rows used in the theorem.}
\label{tab:appendix-local-rows}\\
\toprule
$p$&multiplicity&degree $ef$&$e$&$f$&$d$&$fd$\\
\midrule
\endfirsthead
\toprule
$p$&multiplicity&degree $ef$&$e$&$f$&$d$&$fd$\\
\midrule
\endhead
$3$&1&3&3&1&3&3\\
&1&6&6&1&7&7\\
&2&9&9&1&18&18\\
\midrule
$5$&2&1&1&1&0&0\\
&3&5&5&1&7&7\\
&1&10&10&1&15&15\\
\midrule
$181$&1&3&3&1&2&2\\
&1&6&3&2&2&4\\
&1&18&3&6&2&12\\
\midrule
$997$&1&3&3&1&2&2\\
&1&6&3&2&2&4\\
&1&18&3&6&2&12\\
\midrule
$2346241$&1&3&3&1&2&2\\
&1&6&3&2&2&4\\
&1&18&3&6&2&12\\
\bottomrule
\end{longtable}

The prime-ideal HNF carriers and prime-vector complements are retained in the
arithmetic evidence rather than typeset: they occupy several megabytes and
do not shorten any proof.  Their canonical aggregate guard is
\[
\texttt{ee43a97eb36fdfb21565b339da78b7b2a4e73ef226e34b49223f127a1617ae4a}.
\]

\subsection{\texorpdfstring{All certified \(\theta_{36}\) factor rows}{All certified theta36 factor rows}}

At $3$ and $5$, the factor rows remain identical at precisions
$900,950,1000$, every factor is monic and simple, and the minimum
coefficientwise multiply-back valuations are exactly $900,950,1000$.  The
compressed rows are:

\begin{table}[ht]
\centering
\caption{Wild $\theta_{36}$ factor rows.  The last column is the
polynomial-discriminant valuation of one factor.}
\label{tab:wild-theta-factor-rows}
\begin{tabular}{@{}c c c c c c@{}}
\toprule
$p$&count&factor degree&mod-$p$ exponent&residual degree&
factor poldisc exponent\\
\midrule
$3$&3&3&3&1&11\\
&1&9&9&1&62\\
&1&18&18&1&269\\
\midrule
$5$&1&1&1&1&0\\
&1&5&5&1&27\\
&3&10&10&1&123\\
\bottomrule
\end{tabular}
\end{table}

Thus the global/twice-largest bounds are $886/538$ at $3$ and $746/246$ at
$5$.  All three precisions exceed both numbers.

At each of $181,997,2346241$, the $\theta_{36}$ rows are
\begin{equation}
\begin{array}{c|rrrr}
\text{factor degree}&3&6&9&18\\ \hline
e&3&3&3&3\\
f&1&2&3&6\\
d&2&2&2&2\\
\text{factor poldisc exponent}&2&4&6&12.
\end{array}
\label{eq:tame-theta-rows}
\end{equation}
Their field-different contributions sum to $24$, the global
polynomial-discriminant exponent is $24$, and twice the largest factor
exponent is $24$.  Precision $40$ clears both bounds, and factor
multiplication is verified at $20,30,40$.

For $\delta_{36}$, the same apparent factor degrees at tame precision have
factor polynomial-discriminant exponents $4,20,48,204$.  Hence the
twice-largest bound is $408$, while the global exponent is $840$.  Precision
$40$ clears neither.  No $\delta_{36}$ row, at any prime, is used to prove
or corroborate a local partition.  This negative authority statement is
part of the premise graph, not an observation omitted from the successful
lane.
